\documentclass[12pt]{article}
\usepackage[utf8]{inputenc}
\usepackage[T1]{fontenc}
\usepackage{hyperref}
\title{Northern Toussian Language Resources for K-12 Curriculum Development}
\author{Compiled by Claude Sonnet 5 (claude-sonnet-5) for Andr\'as Kornai}
\date{}
\begin{document}
\maketitle

\section{Introduction}
Northern Toussian (ISO 639-3 tsp; Glottocode nort2787) is spoken in the Hauts-Bassins region of Burkina Faso, around Orodara. It is also called Tusya, Tusian or Win. Its Niger-Congo affiliation is uncertain. Speaker estimates vary widely, from about 18,500 to about 46,000. Children still acquire the language from infancy in most villages. The searches found no school materials, primers, dictionary or standard orthography for the northern variety. This dossier is therefore short.

\section{Primary Educational Resources}
No primers, textbooks or school programmes for Northern Toussian turned up in the searches. None are claimed here.

\section{Linguistic and Reference Materials}
\begin{itemize}
\item Overview article on the Tusya language (Wikipedia): \url{https://en.wikipedia.org/wiki/Tusya_language}
\item ``A Preliminary Account of the Northern Toussian Balafon Surrogate Language'' (Frontiers in Communication, 2021). It documents the griots' musical encoding of speech and gives linguistic background on the language: \url{https://www.frontiersin.org/journals/communication/articles/10.3389/fcomm.2021.652740/pdf}
\item Northern Toussian entry in the numerals database of the Max Planck Institute for Evolutionary Anthropology: \url{https://lingweb.eva.mpg.de/channumerals/Toussian-Northern.htm}
\end{itemize}

\section{Religious and Community Resources}
Joshua Project has a profile of the Northern Toussian people of Burkina Faso, with population estimates: \url{https://joshuaproject.net/people_groups/15573/UV}. No scripture translations or community organisations were found.

\section{Digital and Audio Resources}
No audio, video or corpus resources for Northern Toussian were found. Wikitongues lists the related Southern Toussian (wib), not the northern variety: \url{https://wikitongues.org/languages/wib/}

\section{Curriculum Development Context}
There is little to build on. A curriculum would first need an agreed orthography, a lexicon and basic descriptive work, ideally with the community and Burkinab\'e national-language education bodies. The 2021 balafon study may be a useful starting point for cultural content. Very little published material exists, and field documentation would be needed.
\end{document}